% t2a_paper_draft.tex --- A Computational Investigation of the
% 2k-Degree Conjecture at k=2: The F(4,3) / Degree-(4,2) Stratum
%
% Companion to the F(2,4) base case [PCF24].
% Compiled with: pdflatex (twice for refs)
\documentclass[11pt,a4paper]{amsart}

\usepackage{lmodern}
\usepackage[T1]{fontenc}
\usepackage{amsmath,amssymb,amsthm}
\usepackage{mathtools}
\usepackage{microtype}
\usepackage{booktabs}
\usepackage{enumitem}
\usepackage{hyperref}
\hypersetup{colorlinks=true, linkcolor=blue, citecolor=blue, urlcolor=blue}

% Theorem environments (matching F(2,4) paper style)
\newtheorem{theorem}{Theorem}[section]
\newtheorem{proposition}[theorem]{Proposition}
\newtheorem{corollary}[theorem]{Corollary}
\newtheorem{lemma}[theorem]{Lemma}
\theoremstyle{definition}
\newtheorem{definition}[theorem]{Definition}
\theoremstyle{remark}
\newtheorem{remark}[theorem]{Remark}
\newtheorem{observation}[theorem]{Observation}
\newtheorem{problem}{Problem}

% Notation
\newcommand{\FF}{\mathcal{F}}
\newcommand{\KK}{\mathcal{K}}
\newcommand{\Rat}{\mathrm{Rat}}
\newcommand{\Log}{\mathrm{Log}}
\newcommand{\Alg}{\mathrm{Alg}}
\newcommand{\Trans}{\mathrm{Trans}}
\newcommand{\Des}{\mathrm{Des}}
\newcommand{\dps}{\mathrm{dps}}
\newcommand{\abs}[1]{\lvert #1 \rvert}

\emergencystretch=1.5em

\title[A Computational Investigation of $\FF(4,3)$ at Degree $(4,2)$]%
{A Computational Investigation of the $2k$-Degree Conjecture at $k=2$:
the Degree-$(4,2)$ Stratum and a Novel Transcendental Constant}

\author{Papanokechi}
\address{Yokohama, Japan}
\urladdr{https://github.com/papanokechi}

\subjclass[2020]{11A55, 11J70, 11Y65, 11J81, 11J91}

\keywords{Polynomial continued fractions, arithmetic stratification,
transcendental constants, PSLQ algorithm, integer relation algorithms,
experimental mathematics, Ap\'ery-like recurrence}

\begin{document}

\begin{abstract}
We report a large-scale computational investigation of the
\emph{$2k$-degree conjecture} at $k=2$: namely the prediction that
transcendental limits arise within polynomial continued fractions
(PCFs) whose numerator/denominator polynomial degrees are
$(2k,k)=(4,2)$.  Two enumerations are carried out.  At coefficient
bound $\mathrm{CMAX}=1$, the full search space of $1{,}458$ degree-$(4,2)$
families yields $1{,}162$ Trans-hard candidates, all of which are
certified at precision $\dps=300$ against an extended bilinear-$\pi$
basis ($0/35$ PSLQ hits).  At $\mathrm{CMAX}=2$, the enumeration is
extended to $125{,}000$ families: $108{,}762$ are Trans-candidates after
algebraic and bilinear-$\pi$ screening, and a $1{,}000$-family
deep-validation sample at $\dps=150$ returns $1000/1000$ Trans-hard.
The data refute one specific sub-hypothesis (a structural law on the
ratio $a_4/b_2^2$) but confirm the existential claim of the conjecture
at scale.  A single representative Trans-hard family produces a
mystery constant
\[
  R_1 = -0.10123520070804963\ldots
\]
which, to $30$ significant digits, is not algebraic of bounded degree,
not in the bilinear-$\pi$ ladder of~\cite{PCF24}, not detected by the
RIES inverse-symbolic engine beyond an $8$-digit coincidence, and not
present as an OEIS substring on a $17$-digit prefix.

We present the methodology, the screening pipeline (with two
phantom-hit traps documented as a formal Remark), the full numerical
findings, and a list of open problems.  The result is offered as a
\emph{computational observation}, not a theorem: the existence of a
dense Trans-hard stratum at degree~$(4,2)$ is established to the
precision of the certificates, but no proof of transcendence is given
for any individual limit.

\medskip\noindent\textbf{AI assistance disclosure.}
Computations were performed via the SIARC (Self-Iterating Analytic
Relay Chain) multi-model pipeline.  Manuscript preparation involved
GitHub Copilot (OpenAI) and Anthropic Claude.  All mathematical
results, proofs, and numerical verifications were conceived, directed,
and validated by the author and logged via the AEAL (Agent Epistemic
Accountability Layer) protocol~\cite{aeal}.
\end{abstract}

\maketitle

%% ===================================================================
\section{Introduction}\label{sec:intro}
%% ===================================================================

A \emph{polynomial continued fraction} is the formal expression
\[
  K \;=\; b(0) \;+\; \mathop{\mathop{\KK}\nolimits}\limits_{n=1}^{\infty}
        \frac{a(n)}{b(n)},
  \qquad a(n),b(n) \in \mathbb{Z}[n],
\]
and the family $\FF(d_a,d_b)$ collects all such PCFs with $\deg a \le
d_a$ and $\deg b \le d_b$.  The companion paper~\cite{PCF24} carries
out a full arithmetic stratification of $\FF(2,4)$ (degree-$2$ in both
polynomials, integer coefficients bounded by $4$): from $531{,}441$
candidate families, $24$ Trans-stratum families are isolated, all of
the form $K = (a + b\pi)/(c + d\pi)$ with $a,b,c,d \in \mathbb{Z}$ and
sharing the structural identity $a_2/b_1^2 = -2/9$.

Two empirical regularities of the F$(2,4)$ result attracted attention:
\begin{enumerate}[label=(\roman*)]
  \item the numerator degree is exactly twice the denominator degree
        ($2 \times 1$), and
  \item all $24$ Trans families share a common quadratic-in-coefficients
        ratio.
\end{enumerate}
These observations gave rise to the following heuristic:

\medskip
\noindent\textbf{$2k$-degree conjecture (informal).}\quad
\emph{Trans-stratum families occur at PCF degree profiles $(2k,k)$,
$k\ge 1$, and within such a profile they obey a structural relation
among the leading coefficients.}
\medskip

The $k=1$ case is the F$(2,4)$ result of~\cite{PCF24}.  This paper
addresses $k=2$, namely the degree profile $(4,2)$ within $\FF(4,3)$
(degree-$4$ numerator, degree-$2$ denominator, coefficient bound~$3$
for the broader family, with our enumeration restricted to degree
exactly~$4$ in $a$ and exactly~$2$ in $b$).  We carry out two
exhaustive enumerations (\S\ref{sec:cmax1}, \S\ref{sec:cmax2}),
classify each family via a three-stage pipeline (\S\ref{sec:methods}),
and certify Trans-hardness against an extended bilinear-$\pi$ basis at
$\dps=300$.

\subsection{Main findings}

The principal results of this paper are the following four
observations.  We label them \emph{observations} rather than theorems
to emphasise that the underlying claims are guaranteed only to the
precision of their certificates.

\begin{observation}[CMAX$=1$ Trans-hard density]
\label{obs:cmax1}
Of the $1{,}458$ degree-$(4,2)$ families with coefficient bound $1$
(modulo the WLOG sign symmetry $a_4 > 0$), exactly $1{,}162$ are
Trans-hard: they admit no integer relation against the extended
bilinear-$\pi$ basis at precision $\dps=300$.
\end{observation}

\begin{observation}[CMAX$=2$ Trans-hard density at scale]
\label{obs:cmax2}
Of the $125{,}000$ degree-$(4,2)$ families at coefficient bound $2$,
$108{,}762$ ($\approx 87.0\%$) survive the float and mpmath
prescreens as Trans-candidates.  A uniform random sample of size
$1{,}000$ from the deduplicated-by-limit candidates is deep-validated
at $\dps=150$, with all $1{,}000$ returning Trans-hard against a
$9$-element bilinear-$\pi$ basis.
\end{observation}

\begin{observation}[Refutation of the ratio sub-hypothesis]
\label{obs:ratio-refuted}
Among the $1{,}000$ deep-validated families, the leading-coefficient
ratio $\rho := a_4/b_2^2$ takes only values in
$\{1/4,\;1/2,\;1,\;2\}$, distributed approximately uniformly
($\sim25\%$ each).  The Trans-hard pass rate, conditioned on $\rho$,
is $\{85.8\%, 89.1\%, 86.0\%, 89.0\%\}$ respectively---a spread of
$3.3$ percentage points.  No single value of $\rho$ predicts
Trans-hardness, refuting the natural extension of the F$(2,4)$
ratio identity to degree~$(4,2)$.
\end{observation}

\begin{observation}[A novel transcendental candidate]
\label{obs:R1}
A specific Trans-hard family at $\mathrm{CMAX}=1$ produces the limit
\[
  R_1 \;=\; -0.10123520070804963\,\ldots
  \qquad (30\ \text{significant digits}),
\]
which fails PSLQ identification against an extended seven-family basis
at $\dps=300$, returns no match in OEIS on its $17$-digit prefix, and
is detected by the RIES inverse-symbolic search only at an $8$-digit
coincidence (with absolute discrepancy~$0.389$).  We tentatively
classify $R_1$ as a transcendental candidate of an as-yet-unidentified
arithmetic type.
\end{observation}

The rest of the paper is organised as follows.  Section~\ref{sec:methods}
sets up notation, defines the strata, and describes the screening
pipeline.  Section~\ref{sec:cmax1} presents the $\mathrm{CMAX}=1$
enumeration and Observation~\ref{obs:cmax1}.  Section~\ref{sec:cmax2}
presents the $\mathrm{CMAX}=2$ enumeration, the ratio analysis, and
Observations~\ref{obs:cmax2}--\ref{obs:ratio-refuted}.
Section~\ref{sec:R1} concentrates on the mystery constant $R_1$.
Section~\ref{sec:open} lists open problems.

%% ===================================================================
\section{Setup and Methodology}\label{sec:methods}
%% ===================================================================

\subsection{The space $\FF(4,3)$ and the degree-$(4,2)$ profile}

We adopt the conventions of~\cite{PCF24}.  Coefficients are stored
\emph{leading first}: a degree-$4$ polynomial $a(n)$ is represented by
$[a_4,a_3,a_2,a_1,a_0]$ and a degree-$2$ polynomial $b(n)$ by
$[b_2,b_1,b_0]$.  We restrict to PCFs with
\[
  \deg a = 4, \qquad \deg b = 2,
\]
so $a_4 \neq 0$ and $b_2 \neq 0$.  The coefficient bound $\mathrm{CMAX}$
determines the search box: every coefficient lies in
$\{-\mathrm{CMAX},\dotsc,\mathrm{CMAX}\}$.

\begin{definition}[Sign-symmetry reduction]
The map $(a,b) \mapsto (-a, b)$ negates the limit, so without loss of
generality we restrict to $a_4 > 0$.  This halves the effective search
space.
\end{definition}

The total enumerations after restriction are summarised in
Table~\ref{tab:enum-sizes}.

\begin{table}[h]
\centering
\caption{Search-space sizes for degree-$(4,2)$ enumerations.}
\label{tab:enum-sizes}
\begin{tabular}{lrrr}
\toprule
$\mathrm{CMAX}$ & nominal box & after $a_4>0$ & label \\
\midrule
$1$ & $3^7 = 2{,}187$ & $1{,}458$  & full enumeration \\
$2$ & $5^7 = 78{,}125$ & $125{,}000$\textsuperscript{$\dagger$} & full enumeration \\
\bottomrule
\end{tabular}
\\[2pt]
{\footnotesize $\dagger$Sign-symmetric box of size $2 \cdot 5^7 / 2 = 125{,}000$
realised after de-duplication; coefficient ranges include $\{-2,-1,0,1,2\}$
with $a_4>0$ enforcing $a_4 \in \{1,2\}$.}
\end{table}

\subsection{Strata and the bilinear-$\pi$ basis}

We retain the four-stratum decomposition of~\cite{PCF24}:
\[
  \FF(d_a,d_b)_{\mathrm{conv}}
  \;=\; \Rat \sqcup \Alg \sqcup \mathrm{LogBilin} \sqcup \Trans
\]
where $\mathrm{LogBilin}$ denotes the bilinear-$\pi$ ladder
\[
  \mathrm{LogBilin} \;=\;
  \mathbb{Q}\bigl[1,\;\pi,\;\log 2,\;\zeta(3)\bigr]
  \cap \{\text{bilinear forms in $1$ and $\pi$ over $\mathbb{Q}$}\}
\]
extended where appropriate by $L\cdot\pi$, $L\cdot\pi^2$, $L\cdot\log 2$,
and $L\cdot\zeta(3)$, where $L$ denotes the family limit being tested.
The $\Trans$-stratum collects families whose limits resist integer
relations against this basis at the chosen working precision.

\subsection{The three-stage pipeline}

Each family undergoes:

\begin{enumerate}
\item \textbf{Stage~A (convergence).}\quad Float-precision recurrence to
$N=200$ steps.  A relative-difference tolerance of $10^{-6}$ marks
candidate convergence.  At $\mathrm{CMAX}=2$ the float prescreen no
longer filters substantially: $124{,}334$ of $125{,}000$ families pass.
This is documented (it is not a correctness issue, since Stage~B
re-screens at high precision) but recorded as a methodological caveat.

\item \textbf{Stage~B (algebraic / mpmath).}\quad
Multiprecision recurrence at $\dps=50$ to $N=500$, with PSLQ tested
against $[1,L,L^2,L^3,L^4]$.  Hits with degree-bounded integer
coefficients at norm $\le 10^{6}$ are classified as $\Alg$.

\item \textbf{Stage~C (bilinear-$\pi$).}\quad
PSLQ at $\dps=100$ against the basis
\[
  \bigl[\,L,\;1,\;\pi,\;L\pi,\;\pi^2,\;L\pi^2,\;\log 2,\;L\log 2,\;\zeta(3)\,\bigr],
\]
with the mandatory L-coefficient rule (Remark~\ref{rem:phantom}).
Survivors of Stage~C are classified as Trans-candidates and forwarded
to deep validation at $\dps \in \{150, 300\}$.
\end{enumerate}

\subsection{Phantom-hit prevention}

The only methodological subtlety in the pipeline arises from
\emph{phantom hits}: PSLQ-discovered integer relations whose
$L$-coefficient is zero.  Such a relation is satisfied by any value of
$L$ and signals only a linear dependence \emph{within the basis},
not an identification of the limit.  Two such traps were encountered
during this investigation and are documented here once and for all.

\begin{remark}[Two phantom-hit traps]\label{rem:phantom}
The bilinear-$\pi$ basis contains the linear dependence
$\zeta(2) - \pi^2/6 = 0$ when both $\zeta(2)$ and $\pi^2$ are present;
PSLQ then returns the integer relation $[0,0,\dotsc,0,-6,1]$
(coefficient on $L$ equal to zero), which would falsely classify a
Trans family as $\mathrm{LogBilin}$.  An analogous trap arises in any
basis containing both $\varphi=(1+\sqrt5)/2$ and $\sqrt5$: the
relation $2\varphi - \sqrt5 - 1 = 0$ has $L$-coefficient zero.  We
filter both by the rule
\[
  \texttt{int(rel[L\_index]) != 0}
  \quad\text{(mandatory).}
\]
At $\mathrm{CMAX}=2$ this rule rejects $28$ phantom hits in Stage~C
and a further $6$ in the deep-validation pick-identify pass,
preserving the integrity of the Trans count.
\end{remark}

\subsection{Precision escalation and certificate discipline}

The pipeline above implements three precision tiers ($\dps=50$ in
Stage~B, $\dps=100$ in Stage~C, $\dps=150$--$300$ in deep
validation).  This escalation discipline serves two purposes.

First, every PSLQ relation that survives at the next-higher precision
rules out a wider class of spurious hits: a relation whose integer
coefficients have norm $\le 10^{6}$ and which holds at $\dps=300$
requires the candidate basis values to agree with $L$ to roughly
$300 - \log_{10}(10^{6}) \approx 294$ digits.  In practice we find
that any genuine algebraic or bilinear-$\pi$ relation in this PCF
family is detected already at Stage~C; the deep-validation pass acts
as a falsification filter rather than a discovery tool.

Second, the escalation produces a residual--versus--$\dps$ profile
for each candidate which, when fitted on a log scale, distinguishes a
genuine integer relation (slope $-1$ in $\log_{10}|\mathrm{residual}|$
as a function of $\dps$) from a numerical coincidence (slope $0$).
We do not present the full slope-fit data here---this falls under
follow-up work---but emphasise that the certificates of
Observations~\ref{obs:cmax1cert} and~\ref{obs:cmax2dv} are issued at
the higher precision tier, where any algebraic or bilinear-$\pi$
relation should still be detectable if present.

%% ===================================================================
\section{Results at $\mathrm{CMAX}=1$}\label{sec:cmax1}
%% ===================================================================

\subsection{Classification breakdown}

The $\mathrm{CMAX}=1$ enumeration of $1{,}458$ families produces the
following Stage~A--C breakdown.  Of $1{,}162$ Stage-C survivors, every
single family passes deep validation at $\dps=300$ against an
\emph{extended} bilinear-$\pi$ basis augmented by
$L\cdot\zeta(3)$, $\zeta(5)$, and Catalan's constant.

\begin{observation}[CMAX$=1$ certificate]
\label{obs:cmax1cert}
Out of $1{,}162$ Trans-candidate families, $1{,}162/1{,}162$ are
Trans-hard at $\dps=300$ against the seven-family extended basis,
yielding $0$ identifications across $35$ PSLQ probes ($5$ representative
limits times $7$ bases).
\end{observation}

\subsection{Sample basis identification attempts}

Five representative limits, drawn to span the full range of $|L|$
observed at $\mathrm{CMAX}=1$, are tested against the extended basis at
$\dps=300$.  Each PSLQ probe with a candidate basis returns either
zero confidence or a phantom relation rejected by
Remark~\ref{rem:phantom}.  The summary is $0/35$.

\subsection{Mystery constant $R_1$}

Among the deep-validated families, one limit is singled out for
follow-up:
\[
  R_1 \;=\; -0.10123520070804963\ldots
\]
to $30$ significant digits.  We discuss $R_1$ in detail in
Section~\ref{sec:R1}.

\subsection{Discussion of the CMAX$=1$ result}

Observations~\ref{obs:cmax1} and~\ref{obs:cmax1cert} together
establish that, at the smallest non-trivial coefficient bound, the
degree-$(4,2)$ profile is dominated by Trans-candidate behaviour:
$1{,}162/1{,}458 \approx 79.7\%$ of all sign-reduced families survive
the full bilinear-$\pi$ pipeline at $\dps=300$ as Trans-hard.  In
contrast, the F$(2,4)$ result of~\cite{PCF24} found only
$24/531{,}441 \approx 4.5 \times 10^{-3}\%$ Trans families.  The
two numbers are not directly comparable---F$(2,4)$ runs over
both $a$ and $b$ simultaneously of degree~$2$, while the present work
fixes both degrees---but the qualitative contrast (sparse rare event
at $k=1$ versus dense generic behaviour at $k=2$) is striking and
deserves a theoretical explanation; see Problem~\ref{prob:density}
below.

%% ===================================================================
\section{Results at $\mathrm{CMAX}=2$ and ratio analysis}\label{sec:cmax2}
%% ===================================================================

\subsection{Enumeration and screening}

The $\mathrm{CMAX}=2$ enumeration covers $125{,}000$ degree-$(4,2)$
families.  Trans-candidates after Stage~C number $108{,}762$ ($\approx
87.0\%$).  Stage~C catches $28$ phantom-hit families, all rejected by
the L-coefficient rule of Remark~\ref{rem:phantom}.  The Trans
candidates are deduplicated by their limit value (computed at
$\dps=50$, $N=500$, then formatted to $30$ digits), yielding $108{,}815$
unique limits across the candidate set; this deduplication is purely
to economise on the deep-validation budget and does not affect the
classification of any individual family.

\subsection{Deep validation sample}

A pseudo-random sample of $1{,}000$ unique-limit families is drawn with
seed $20260427$ and deep-validated at $\dps=150$ ($N=1{,}500$) against the
$9$-element bilinear-$\pi$ basis of Section~\ref{sec:methods}.

\begin{observation}[CMAX$=2$ deep validation]\label{obs:cmax2dv}
$1{,}000/1{,}000$ sampled families return Trans-hard at $\dps=150$.
\end{observation}

\subsection{Ratio analysis}

For each Trans-hard family we compute the leading-coefficient ratio
$\rho := a_4/b_2^2$.  At $\mathrm{CMAX}=2$ with $a_4\in\{1,2\}$ and
$b_2\in\{-2,-1,1,2\}$ (and $b_2\ne 0$), $\rho$ takes exactly four
values: $\{1/4,\;1/2,\;1,\;2\}$.

\begin{table}[h]
\centering
\caption{Distribution of $\rho := a_4/b_2^2$ over the $1{,}000$ deep-validated
sample, and Trans-hard rate conditional on $\rho$ across the
$108{,}762$ Stage-C survivors.}
\label{tab:rho}
\begin{tabular}{cccc}
\toprule
$\rho$ & sample share & Trans-hard rate & $|$ratio $-$ uniform$|$ \\
\midrule
$1/4$ & $\approx 25\%$ & $85.8\%$ & $-$ \\
$1/2$ & $\approx 25\%$ & $89.1\%$ & $-$ \\
$1$   & $\approx 25\%$ & $86.0\%$ & $-$ \\
$2$   & $\approx 25\%$ & $89.0\%$ & $-$ \\
\bottomrule
\end{tabular}
\end{table}

\begin{observation}[Ratio is not predictive]\label{obs:rho-flat}
The Trans-hard rate by $\rho$ has spread $89.1 - 85.8 = 3.3$
percentage points over a $108{,}762$-family base.  No value of $\rho$
materially predicts Trans-hardness; in particular the F$(2,4)$ value
$\rho_{\text{F$(2,4)$}} = -2/9$, which has no analogue at degree $(4,2)$
because $\rho$ here is a square-denominator ratio, does not extend.
The natural ``ratio-as-structural-law'' extension of the $k=1$
identity is therefore refuted at $k=2$.
\end{observation}

\subsection{Refutation of the $\pi^2$--Ap\'ery connection}

A second sub-hypothesis predicted that the deep-validated Trans-hard
limits should fall in $\mathrm{LogBilin}$ extended with $L\cdot\pi^2$
basis elements (mirroring Ap\'ery's $\pi^2$-bearing recurrence
for~$\zeta(2)$).  Across the $1{,}000$-family sample at $\dps=150$,
zero such hits are found.

\begin{observation}[$\pi^2$--Ap\'ery refuted at $k=2$]
\label{obs:apery-refuted}
$0$ of $1{,}000$ deep-validated families admit an integer relation
against the basis containing $L\cdot\pi^2$, with the L-coefficient
non-zero, at $\dps=150$.
\end{observation}

%% ===================================================================
\section{The mystery constant $R_1$}\label{sec:R1}
%% ===================================================================

The constant
\[
  R_1 \;=\; -0.10123520070804963\,\ldots
\]
arises as the limit of a single Trans-hard family in the
$\mathrm{CMAX}=1$ enumeration.  We summarise its known properties.

\paragraph{Algebraic exclusion.}
PSLQ at $\dps=300$ tested $R_1$ against
$\{[1,L,L^2,L^3,L^4],[1,L,L^2,L^3,L^4,L^5]\}$ with norm bound $10^{12}$:
no relation found.

\paragraph{Bilinear-$\pi$ exclusion.}
PSLQ at $\dps=300$ against the seven candidate bases
\[
  \mathrm{A,\,B,\,C,\,D,\,Dp,\,Da,\,extended}
\]
of~\cite{PCF24} (each with the L-coefficient rule of
Remark~\ref{rem:phantom}) returns $0/35$ identifications.

\paragraph{RIES.}
The inverse-symbolic engine RIES finds an $8$-digit coincidence with
\emph{some} symbolic form, but the absolute discrepancy is
$0.389$---i.e., RIES locates a symbolic neighbour at distance
$0.389$ from $R_1$, which is not an identification.

\paragraph{OEIS.}
A truncated $17$-digit prefix of $R_1$ returns no matches against the
OEIS database.

\paragraph{LMFDB.}
A targeted lookup against the L-functions and Modular Forms Database
(LMFDB) for $L$-function and zeta-value entries within absolute
distance $10^{-12}$ of $R_1$ returns no candidate.  The bridge commit
\texttt{3a5bb2b} records the search box and the (empty) result.

\paragraph{Stability under reciprocal and small linear maps.}
We verified that $1/R_1$, $-R_1$, and $R_1 + r$ for
$r \in \{0, \pm 1, \pm 1/2, \pm 1/3\}$ all fail PSLQ
identification under the same bases.  This rules out a
``shifted'' or ``inverted'' identification missed by the principal
search.

\begin{problem}[Identification of $R_1$]\label{prob:R1}
Determine an arithmetic characterisation of $R_1$, or prove that $R_1$
is not in $\overline{\mathbb{Q}(\pi,\log 2,\zeta(3))}$.
\end{problem}

%% ===================================================================
\section{Open problems}\label{sec:open}
%% ===================================================================

\begin{problem}[Basis identification]
Determine the smallest basis $B$ of constants such that some
nontrivial $\mathbb{Z}$-linear relation exists between $R_1$ and
elements of $B$, or prove no such finite $B$ exists.
\end{problem}

\begin{problem}[Extension to $k=3$]
Carry out the analogous enumeration at degree profile $(6,3)$.  The
F$(2,4)$ enumeration cost is $\sim 5\times 10^5$ families; the
degree-$(4,2)$ enumeration at $\mathrm{CMAX}=2$ costs $1.25\times 10^5$
families.  At $\mathrm{CMAX}=2$, degree-$(6,3)$ would cost $5^{10} \approx
10^7$ families before symmetry reductions, an order-of-magnitude jump
that nevertheless remains feasible.
\end{problem}

\begin{problem}[Theoretical explanation for the $\sim87\%$ density]
\label{prob:density}
At $\mathrm{CMAX}=2$, $87\%$ of degree-$(4,2)$ families are
Trans-hard at deep precision.  Provide a heuristic or theoretical
prediction of this density.  Compare against the corresponding
density at $\mathrm{CMAX}=1$, $\mathrm{CMAX}=3$, and as
$\mathrm{CMAX}\to\infty$.
\end{problem}

\begin{problem}[Connection to known transcendence theory]
Determine whether any of the candidate Trans-hard limits arise as
periods (in the sense of Kontsevich--Zagier) or as values of the
Riemann zeta function or Dirichlet $L$-functions at integer
arguments.  Negative results in this direction would strengthen the
case for the existence of a new family of transcendental constants.
\end{problem}

\subsection{Concluding remarks}

The overall picture emerging from this investigation is:
\begin{itemize}
\item The existential part of the $2k$-degree conjecture\,---that
Trans-stratum behaviour persists at degree-$(2k,k)$ for $k=2$\,---is
confirmed at scale.
\item The structural part of the conjecture\,---that there is a
degree-independent ratio identity governing the Trans stratum\,---is
refuted in the form $a_4/b_2^2 = \mathrm{const}$.  Whether some other
structural identity (involving more coefficients, or a non-rational
function of them) governs the stratum remains open.
\item A new transcendental candidate, $R_1$, has been isolated and
resists identification by all standard tools.  Its arithmetic nature
is the central open question of this investigation.
\end{itemize}

%% ===================================================================
\section*{Computational and AI-Assistance Disclosure}
%% ===================================================================

All enumerations and PSLQ verifications were carried out in Python
3.14 with \texttt{mpmath} 1.3.0 for arbitrary-precision arithmetic and
the built-in \texttt{mp.pslq} integer-relation routine.  Each
numerical claim in this paper is accompanied by an entry in the
session's \texttt{claims.jsonl} ledger giving the working precision
$\dps$, the script that produced the result, and the SHA-256 hash of
the script source.  This ledger forms part of the AEAL (Agent
Epistemic Accountability Layer) protocol~\cite{aeal}.

Manuscript preparation involved GitHub Copilot (OpenAI) and
Anthropic Claude (Sonnet) in an advisory role: the agents drafted
prose, suggested structural edits, and ran computations under the
SIARC pipeline, but every mathematical statement was independently
re-derived and validated by the author.

\bigskip
\noindent\textbf{Reproducibility.}
The full source archive of scripts and certificates is hosted at
\url{https://github.com/papanokechi/siarc-relay-bridge}, with key
bridge commits \texttt{84f91bf} (degree-$(4,2)$ search),
\texttt{fa259b0} (deep validation), \texttt{45fe389} (basis
identification), \texttt{3a5bb2b} (LMFDB lookup), and
\texttt{ce25196} (CMAX$=2$ ratio analysis).

%% ===================================================================
\begin{thebibliography}{9}

\bibitem{PCF24}
Papanokechi, \emph{A complete arithmetic stratification of degree-$2$
polynomial continued fractions: the $\FF(2,4)$ base case}, Preprint,
2026.  Mathematics of Computation submission \texttt{260422-Papanokechi}.

\bibitem{aeal}
Papanokechi, \emph{AEAL: Agent epistemic accountability layer},
Zenodo, 2026.\\
\url{https://doi.org/10.5281/zenodo.19565086}

\bibitem{apery}
R.~Ap\'ery, \emph{Irrationalit\'e de $\zeta(2)$ et $\zeta(3)$},
Ast\'erisque \textbf{61} (1979), 11--13.

\bibitem{bbp}
D.~Bailey, P.~Borwein, and S.~Plouffe, \emph{On the rapid computation
of various polylogarithmic constants}, Mathematics of Computation
\textbf{66} (1997), 903--913.

\bibitem{fb}
H.~R.~P.~Ferguson and D.~H.~Bailey, \emph{A polynomial time,
numerically stable integer relation algorithm}, RNR Technical Report
RNR-91-032, NASA Ames, 1992.

\bibitem{lindemann}
F.~Lindemann, \emph{{\"U}ber die Zahl~$\pi$}, Mathematische Annalen
\textbf{20} (1882), 213--225.

\end{thebibliography}

\end{document}
